\section{引言}\label{sec:intro}

\subsection{研究背景与统计问题}\label{sec:intro-bg}

时间序列的结构变化既可能表现为均值、方差等低阶统计量的改变，也可能表现为局部轨迹几何、周期结构或状态空间形状的改变。滑动窗口和时延嵌入为后一类信号提供了局部化表示：每个时间片段被转换为点云，持续同调进一步从多尺度上提取连通分支和环等结构信息。相关理论奠定了滑动窗口持久性分析的基础\cite{PereaHarer2015}，持续图像等有限维表示则使持续图能够进入统计扫描和机器学习程序\cite{Adams2017}。近年来，拓扑摘要已被用于离线变点检测、在线监测及依赖序列中的结构异常识别\cite{IslambekovYuvarajGel2020,ZhengMakXieXie2023,ChazalLevrardRoyer2024,NieYue2026}。

为提高候选变点的扫描密度，实际分析常取步长 $s$ 小于窗口宽度 $W$。此时，相邻窗口共享原始观测，窗口特征序列由此形成局部依赖。设第 $i$ 个窗口经固定映射得到特征 $Y_i$。对全部窗口顺序实施置换时，有限样本精确性通常依赖联合分布对排列群保持不变，即
\begin{equation}\label{eq:intro-exchangeability}
 (Y_1,\ldots,Y_M)\overset{d}{=}(Y_{\pi(1)},\ldots,Y_{\pi(M)}),
 \qquad \forall\,\pi\in\mathfrak S_M.
\end{equation}
独立同分布是满足该条件的一种充分情形，排列操作直接要求的条件则是窗口特征层面的联合交换性。窗口共享观测可诱发局部依赖，但特定映射是否具有交换性取决于其联合分布；平稳性、边际同分布以及原始序列的独立同分布性质，在窗口重叠时均不足以一般性地推出式\eqref{eq:intro-exchangeability}。由此形成本文关注的统计错配：科学零假设定义在原始过程的“无变化”层面，置换校准所需的对称性定义在派生窗口特征的联合分布层面。依赖时间序列中的置换理论同样表明，未经修正的排列参考分布可能无法维持名义检验水平\cite{RomanoTirlea2022}；随机排列检验的有限样本有效性还需相应的群不变性条件\cite{HemerikGoeman2018}。

本文考察一种基于持续同调特征的变点检验：每个局部时间窗口映射为点云，点云经持续同调及固定向量化形成特征序列，再通过标准化部分和的最大值检测前后特征均值差异。持续同调决定检验保留的结构信息，重采样程序决定统计量的参考分布。本文重点评价窗口重叠对后者有限样本校准的影响，并以均值--协方差特征作为非拓扑对照。

\subsection{相关工作与研究缺口}\label{sec:intro-lit}

与本文相关的研究可概括为三条脉络。第一类工作研究滑动窗口嵌入、持续同调及持续图向量化，为局部时间片段的拓扑表示提供方法基础\cite{PereaHarer2015,Adams2017}，其关注点主要在于表示的稳定性与可解释性，而未系统考察重叠窗口下推断有效性的有限样本校准。第二类工作将拓扑摘要引入离线或在线变化检测，并显式处理时间依赖与误报控制：文献\cite{NieYue2026}已经分析重叠窗口下的依赖继承，并利用移动块自助法校准在线监测控制限；其它近邻工作则聚焦表示对结构变化的辨识能力及特定检测程序的性能\cite{IslambekovYuvarajGel2020,ZhengMakXieXie2023,ChazalLevrardRoyer2024}。本文沿用这一研究背景，固定表示、离线最大扫描和校准规则，比较两种观测预算下的零假设拒绝率与经验分类。第三类工作研究依赖数据下的置换、dependent wild bootstrap 与依赖乘子自助法，说明时间依赖会改变随机化参考分布，并为若干均值、光滑函数或经验过程问题建立渐近理论\cite{RomanoTirlea2022,Shao2010,BucherKojadinovic2016}。这些理论为窗口特征序列提供了渐近分析的起点，但其适用性仍需针对本文所用的局部映射、扫描泛函与带宽规则逐项验证。

这些结果尚未提供本文实现的具体有限样本保证：重叠局部窗口依次经过点云映射、持续同调向量化和最大扫描后，朴素置换会产生多大程度的有限样本尺寸失真；该失真如何随重叠程度、原序列依赖、观测预算和表示类型变化；依赖感知校准能够在哪些条件下减弱失真，又会在哪些条件下表现为保守或残留反保守。本文不以非线性变换本身主张该管线脱离既有理论框架，也不从渐近一致性推定当前样本规模下的名义尺寸控制，而以固定实现和预设设计网格上的有限样本校准偏差及其条件变化为考察重点。本文将条件组合 $\theta$ 下的有效检验水平记为 $\alpha_{\mathrm{eff}}(\theta)$，并将映射 $\theta\mapsto\alpha_{\mathrm{eff}}(\theta)$ 称为预设离散设计网格上的条件化有限样本校准图谱。本文不估计连续参数空间中的解析边界；“经验校准状态转换”仅指有序设计维度（主要为步长或重叠率）上相邻格点发生校准分类变化的位置。该图谱属于有限样本、点态和设计域内的刻画，不提供全参数空间的统一保证。

\subsection{研究目标与贡献}\label{sec:intro-contrib}

本文的主要贡献是对固定持续同调变点检验开展条件化有限样本校准评价：在两种预算设计和 28 个唯一零假设配置上，报告朴素置换与依赖乘子的拒绝率、区间及经验分类，并量化两种预算下重叠与非重叠条件之间的预设对比。机制分析通过一般局部映射的依赖继承和窗口均值反例说明交换性风险；备择分析通过低阶矩匹配的形状构造与整体平移负对照，解释固定表示的信号敏感性及不变性。理论推导、经验评价与构造对照分别承担机制、实际表现和解释功能。

全文其余部分安排如下：第~\ref{sec:method} 节定义重叠窗口特征、两类表示、最大扫描统计量及校准方法；第~\ref{sec:theory} 节分析有限依赖、交换性失效与方差机制；第~\ref{sec:protocol} 节说明评价对象和实验设计；第~\ref{sec:h0} 节报告零假设下的有限样本校准结果；第~\ref{sec:h1} 节分析备择下的检出与定位性能；第~\ref{sec:discussion} 节讨论方法含义、适用范围与局限；第~\ref{sec:conclusion} 节总结全文。
